\beginchapter Chapter 26. Summary of\\Math Mode

^^{math mode}
To conclude the survey that was begun in Chapter 24, let us investigate
exactly what \TeX's digestive processes can do when \TeX\ is building
lists in math mode or in display math mode.

\ninepoint
\def\\{\smallbreak\textindent{$\bull$}}
\medbreak
\centerline{$*\qquad*\qquad*$}
\medskip\noindent
Three asterisks, just like those that appear here, can be found near the
end of Chapter~24.
Everything preceding the three asterisks in that chapter applies to
math mode as well as to vertical mode, so we need not repeat
all those rules. In particular, Chapter~24 explains assignment commands,
and it tells how kerns, penalties, marks, insertions, adjustments,
and ``whatsits'' are put into math lists. Our present goal
is to consider the commands that have an intrinsically mathematical
flavor, in the sense that they behave differently in math
mode than they do in vertical or horizontal modes.

Math lists are somewhat different from \TeX's other lists because they
contain three-pronged ``^{atoms}'' (see Chapter~17). Atoms come in thirteen
flavors: Ord, Op, Bin, Rel, Open, Close, Punct, Inner, Over, Under, Acc,
Rad, and Vcent. Each atom contains three ``^{fields}'' called its ^{nucleus},
^{superscript}, and ^{subscript}; and each field is either empty or is filled
with a math symbol, a box, or a subsidiary math list. Math symbols, in turn,
have two components: a family number and a position number.

It's convenient to introduce a few more rules of syntax, in order to specify
what goes into a math list:
\beginsyntax
<character>\is<letter>\alt<otherchar>\alt^|\char|<8-bit number>\alt%
    <chardef token>
<math character>\is^|\mathchar|<15-bit number>\alt<mathchardef token>
  \alt|\delimiter|<27-bit number>
<math symbol>\is<character>\alt<math character>
<math field>\is<filler><math symbol>\alt<filler>|{|<math mode material>|}|
<delim>\is<filler>^|\delimiter|<27-bit number>
  \alt<filler><letter>\alt<filler><otherchar>
\endsyntax
We have already seen the concept of \<character> in Chapter~25.
Indeed, characters are \TeX's staple food: The vast majority of all
commands that reach \TeX's digestive processes in horizontal mode are instances
of the \<character> command, which specifies a number between 0 and~255 that
causes \TeX\ to typeset the corresponding character in the current font.
When \TeX\ is in math mode or display math mode, a \<character> command
takes on added significance: It specifies a number between 0 and~$32767=
2^{15}-1$. This is done by replacing the character number by its
^|\mathcode| value. If the
|\mathcode| value turns out to be $32768=\null$\hex{8000}, however,
the \<character>
is replaced by an ^{active character} token having the original character
code (0 to~255); \TeX\ forgets the original \<character> and expands this
active character according to the rules of Chapter~20.

A \<math character> defines a 15-bit number either by specifying it
directly with ^|\mathchar| or in a previous ^|\mathchardef|, or by
specifying a 27-bit |\delimiter| value; in the latter case, the least
significant 12~bits are discarded.

It follows that every \<math symbol>, as defined by the syntax above,
specifies a 15-bit number, i.e., a number between 0 and~32767.
Such a number can be represented in the form $4096c+256f+a$, where
$0\le c<8$, $0\le f<16$, and $0\le a<256$. If $c=7$, \TeX\ changes $c$~to~0;
and in this case if the current
value of\/ ^|\fam| is between 0 and~15, \TeX\ also replaces $f$~by~|\fam|.
This procedure yields, in all cases, a class number~$c$ between 0~and~6,
a family number~$f$ between 0~and~15, and a position number~$a$ between
0 and~255. \ (\TeX\ initializes the value of\/ |\fam| by implicitly
putting the assignment `|\fam=-1|' at the very beginning of\/ |\everymath|
and |\everydisplay|. Thus, the substitution of\/ |\fam| for~$f$ will occur
only if the user has explicitly changed~|\fam| within the formula.)

A \<math field> is used to specify the nucleus, superscript, or subscript
of an atom. When a \<math field> is a \<math symbol>, the $f$ and~$a$
numbers of that symbol go into the atomic field. Otherwise the \<math field>
begins with a~`|{|', which causes \TeX\ to enter a new level
of grouping and to begin a new math list; the ensuing \<math mode material>
is terminated by a~`|}|', at which point the group ends and the resulting
math list goes into the atomic field. If the math list turns out to be
simply a single Ord atom without subscripts or superscripts,
or an Acc whose nucleus is an Ord, the
enclosing braces are effectively removed.

A \<delim> is used to define both a ``small character'' $a$ in family~$f$
and a ``large character''~$b$ in family~$g$, where $0\le a,b\le255$
and $0\le f,g\le15$; these character codes are used to construct variable-size
^{delimiters}, as explained in Appendix~G\null. If the \<delim> is given
explicitly in terms of a 27-bit number, the desired codes are obtained
by interpreting that number as
$c\cdot2^{24}+f\cdot2^{20}+a\cdot2^{12}+g\cdot2^8+b$, ignoring the value
of~$c$. Otherwise the delimiter is specified as a \<letter> or
\<otherchar> token, and the 24-bit ^|\delcode| value of that character is
interpreted as $f\cdot2^{20}+a\cdot2^{12}+g\cdot2^8+b$.

\smallskip
Now let's study the individual commands as \TeX\ obeys them in math mode,
considering first the ones that have analogs in vertical and/or horizontal
mode:

\\^|\hskip|\<glue>, ^|\hfil|, ^|\hfill|, ^|\hss|, ^|\hfilneg|,
^|\mskip|\<muglue>.\enskip
A glue item is appended to the current math list.

\\\<leaders>\<box or rule>\<horizontal skip>.\enskip
Here ^\<horizontal skip> refers to one of the first five glue-appending
commands just mentioned; the formal syntax for \<leaders> and
for \<box or rule> is given in Chapter~24. A glue item that produces
^{leaders} is appended to the current list.

\\^|\nonscript|. A special glue item of width zero is appended; it will
have the effect of cancelling the following item on the list, if that item
is glue and if the |\nonscript| is eventually typeset in ``script style''
or in ``scriptscript style.''

\\^|\noboundary|. This command is redundant and therefore has no
effect; boundary ligatures are automatically disabled in math modes.

\\^\<space token>.\enskip
Spaces have no effect in math modes.

\\|\|\].\enskip ^^{control space}
A control-space command appends glue to the current list, using the same amount
that a \<space token> inserts in horizontal mode when the space factor is 1000.

\\^\<box>.\enskip
The box is constructed, and if the result is void nothing happens.
Otherwise a new Ord atom is appended to the current math list, and the
box becomes its nucleus.

\\^|\raise|\<dimen>\<box>, ^|\lower|\<dimen>\<box>.\enskip
This acts just like an unadorned \<box> command, except that the new box
being put into the nucleus is also shifted up or down by the specified amount.

\\^|\vcenter|\stretch\<box specification>\stretch
|{|\<vertical mode material>|}|.\enskip
A vbox is formed as if `|\vcenter|' had been `|\vbox|'. Then a new
^{Vcent} atom is appended to the current math list, and the box becomes
its nucleus.

\\\<vertical rule>.\enskip
A ^{rule} is appended to the current list (not as an atom).

\\^|\halign|\<box specification>|{|\<alignment material>|}|.\enskip
This command is allowed only in display math mode, and only when the current
math list is empty. The alignment is carried out exactly as if it were done in
the enclosing vertical mode (see Chapter~24), except that the lines are shifted
right by the ^|\displayindent|. The closing `|}|' may be followed by
optional \<assignment> commands other than ^|\setbox|,
after which `|$$|'~must conclude
the display. \TeX\ will insert the ^|\abovedisplayskip| and
^|\belowdisplayskip| glue before and after the result of the alignment.

\\^|\indent|.\enskip
An empty box of width ^|\parindent| is appended to the current list, as the
nucleus of a new Ord atom.

\\^|\noindent|.\enskip
This command has no effect in math modes.

\\|{|\<math mode material>|}|.\enskip
A character token of category 1, or a control sequence like~|\bgroup|
that has been |\let| equal to such a character token, causes \TeX\ to
start a new level of ^{grouping} and also to begin work on a new
math list. When such a group ends---with `|}|'---\TeX\
uses the resulting math list as the nucleus of a new Ord atom that
is appended to the current list. If the resulting math list is a
single Acc atom, however (i.e., an accented quantity),
that atom itself is appended.

\\\<math symbol>.\enskip
(This is the most common command in math mode; see the syntax near the
beginning of this chapter.) \ A math symbol determines three values,
$c$,~$f$, and~$a$, as explained earlier. \TeX\ appends an atom to the current
list, where the atom is of type Ord, Op, Bin, Rel, Open, Close, or Punct,
according as the value of~$c$ is 0,~1, 2, 3, 4, 5, or~6. The nucleus of
this atom is the math symbol defined by $f$ and~$a$.

\\^\<math atom>\<math field>.\enskip
A \<math atom> command is any of the following:
\begindisplay
|\mathord|\alt|\mathop|\alt|\mathbin|\alt|\mathrel|\alt|\mathopen|\cr
\qquad\alt|\mathclose|\alt|\mathpunct|\alt|\mathinner|%
  \alt|\underline|\alt|\overline|\cr
\enddisplay
\TeX\ processes the \<math field>, then appends a new atom of the
specified type to the current list; the nucleus of this atom contains
the specified field.

\\^|\mathaccent|\<15-bit number>\<math field>.\enskip
\TeX\ converts the \<15-bit number> into $c$, $f$, and~$a$ as it does
with any |\mathchar|. Then it processes the \<math field> and appends a
new Acc atom to the current list. The nucleus of this atom contains
the specified field; the accent character in this atom contains $(a,f)$.

\\^|\radical|\<27-bit number>\<math field>.\enskip
\TeX\ converts the \<27-bit number> into $a$, $f$, $b$, and~$g$ as it does
with any |\delimiter|. Then it processes the \<math field> and appends a
new Rad atom to the current list. The nucleus of this atom contains
the specified field; the delimiter field in this atom contains
$(a,f)$ and $(b,g)$.

\\^\<superscript>\<math field>.\enskip
A \<superscript> command is an explicit or implicit character token of
category~7.  If the current list does not end with an atom, a new Ord atom
with all fields empty is appended; thus the current list will end with an
atom, in~all cases.  The superscript field of this atom should be empty; it
is made nonempty by changing it to the result of the specified \<math field>.

\\^\<subscript>\<math field>.\enskip
A \<subscript> command is an explicit or implicit character token of
category~8.  It acts just like a \<superscript> command, except, of course,
that it affects the subscript field instead of the superscript field.

\\^|\displaylimits|, ^|\limits|, ^|\nolimits|.\enskip These commands are
allowed only if the current list ends with an Op atom. They modify a
special field in that Op atom, specifying what conventions should be used
with respect to limits.  The normal value of that field is |\displaylimits|.

\\^|\/|.\enskip
A kern of width zero is appended to the current list. \ (This will have the
effect of adding the italic correction to the previous character, if the
italic correction wouldn't normally have been added.)

\\^|\discretionary|\<disc text>\<disc text>\<disc text>.
This command has the same effect as in horizontal mode (see Chapter~25), but the
third ^{\<disc text>} must produce an empty list.

\\^|\-|.\enskip
This command is usually equivalent to `|\discretionary{-}{}{}|';
the `|-|' is therefore interpreted as a ^{hyphen}, not as a minus sign.
\ (See Appendix~H.)

\def\s{\hskip0pt plus1pt}
\\^|\mathchoice|\s
$\langle$filler$\rangle$\s|{|\s$\langle$math mode material$\rangle$\s|}|\s
$\langle$filler$\rangle$\s|{|\s$\langle$math mode material$\rangle$\s|}|\break
$\langle$filler$\rangle$|{|$\langle$math mode material$\rangle$|}|
$\langle$filler$\rangle$|{|$\langle$math mode material$\rangle$|}|.
Four math lists, which are defined as in the
second alternative of a \<math field>, are
recorded in a ``choice item'' that is appended to the current list.

\\^|\displaystyle|, ^|\textstyle|, ^|\scriptstyle|,
^|\scriptscriptstyle|.\enskip A style-change item that corresponds to the
specified style is appended to the current list.

\\^|\left|\<delim>\<math mode material>^|\right|\<delim>.\enskip
\TeX\ begins a new group, and processes the \<math mode material>
by starting out with a new math list that begins with a left boundary
item containing the first delimiter. This group must be terminated by
`|\right|', at which time the internal math list is completed with a
right boundary item containing the second delimiter. Then \TeX\
appends an Inner atom to the current list; the nucleus of this atom
contains the internal math list.

\\\<generalized fraction command>.\enskip
This command takes one of six forms:
\begindisplay
^|\over|\alt^|\atop|\alt^|\above|\<dimen>\cr
\qquad\alt^|\overwithdelims|\<delim>\<delim>\cr
\qquad\alt^|\atopwithdelims|\<delim>\<delim>\cr
\qquad\alt^|\abovewithdelims|\<delim>\<delim>\<dimen>\cr
\enddisplay
(See Chapter 17.) \ When \TeX\ sees a \<generalized fraction command> it
takes the entire current list and puts it into the numerator field of a
generalized fraction item. The denominator field of this new item is
temporarily empty; the left and right delimiter fields are set equal
to the specified delimiter codes. \TeX\ saves this generalized fraction
item in a special place associated with the current level of math mode
processing. \ (There should be no other generalized fraction item in
that special place, because constructions like `|a\over b\over c|' are
illegal.) \ Then \TeX\ makes the current list empty and continues to
process commands in math mode. Later on, when the current level of math
mode is completed (either by coming to a~`|$|' or a~`|}|' or a~|\right|,
depending on the nature of the current group), the current list will be
moved into the denominator field of the generalized fraction item that
was saved; then that item, all by itself, will take the place of the
entire list. However, in the special case that the current list began
with |\left| and will end with |\right|, the boundary items will be
extracted from the numerator and denominator of the generalized fraction,
and the final list will consist of three items: left boundary,
generalized fraction, right boundary. \ (If you want to watch the process
by which math lists are built, you might find it helpful to type `^|\showlists|'
while \TeX\ is processing the denominator of a generalized fraction.)

\\^\<eqno>\<math mode material>|$|.\enskip
Here \<eqno> stands for either ^|\eqno| or ^|\leqno|; these commands
are allowed only in display math mode. Upon reading \<eqno>, \TeX\
enters a new level of grouping, inserts the ^|\everymath| tokens,
and enters non-display math mode to put the \<math mode material>
into a math list. When that math list is completed, \TeX\ converts it
to a horizontal list and puts the result into a box that will be used
as the equation number of the current display. The closing |$|~token
will be put back into the input, where it will terminate the display.

\\|$|.\enskip
If \TeX\ is in display math mode, it reads one more token, which must
also be~|$|. In either case, the math-shift command terminates the current
level of math mode processing and ends the current group, which should
have begun with either~|$| or~\<eqno>.  Once the math list is finished, it
is converted into a horizontal list as explained in Appendix~G\null. \TeX\
scans \<one optional space> after completing a displayed formula; this is
usually the implicit space at the end of a line in the input file.

\\^|\unhbox|\<8-bit number>, ^|\unhcopy|\<8-bit number>.\enskip
The specified box register must be void. Nothing happens.

\\None of the above: If any other primitive command of \TeX\ occurs in
math mode, an error message will be given, and \TeX\ will try to recover
in a reasonable way. For example, if a |\par| command appears, or if any
other inherently non-mathematical command is given, \TeX\ will try to
insert a `|$|' just before the offending token; this will lead out of math
mode. On the other hand if a totally misplaced token like |\endcsname| or
|\omit| or |#| appears in math mode, \TeX\ will simply ignore it, after
reporting the error. You might enjoy trying to type some really stupid
input, just to see what happens. \ (Say `|\tracingall|' first, as
explained in Chapter~27, in order to get maximum information.)

\ddangerexercise ^{Powers of ten}: The whole \TeX\ language has now been
summarized completely. To~demonstrate how much you know, name all of the ways
you can think of in which the numbers 10, ^^{Derek, Bo} 100, 1000, 10000,
and 100000 have special significance to \TeX.
\answer Radix 10 notation is used for numeric constants and for the output
of numeric data. The first 10 |\count| registers are displayed at each
|\shipout|, and their values are recorded on the |dvi| file at such times.
% Also, TeX was first implemented on a DEC-10.
% The \catcode for <space> is 10.
% My birthday is January 10.
% Can this all be just a coincidence?
A box whose glue has stretched or shrunk to its stated stretchability
or shrinkability has badness 100; this badness value separates ``loose''
boxes from ``very loose'' or ``underfull'' ones. \TeX\ will scroll up
to 100 errors in a single paragraph before giving up (see Chapter~27).
The normal values of\/ |\spacefactor| and |\mag| are 1000. A~|\prevdepth|
value of $-1000\pt$ suppresses interline glue. The badness rating of a
box is at most 10000, except that the |\badness| of
an overfull box is 1000000.  |INITEX| initializes |\tolerance| to
10000, thereby making all line breaks feasible. Penalties of 10000 or more
prohibit breaks; penalties of $-10000$ or less make breaks mandatory. The
cost of a page break is 100000, if the badness is 10000 and if the
associated penalties are less than 10000 in magnitude (see Chapter~15).

\ddangerexercise ^{Powers of two}: Name all of the ways
you can think of in which the numbers 8, 16, 32, 64, 128, 256, \dots\
have special significance to \TeX.
\answer \TeX\ allows constants to be expressed in radix 8 (octal)
or radix~16 (hexadecimal) notation, and it uses hexadecimal notation to
display |\char| and |\mathchar| codes. There are 16 families for math
fonts, 16 input streams for |\read|, 16 output streams for |\write|.
A |\catcode| value must be less than~16. The notation |^^?|, |^^@|,
|^^A| specifies characters whose codes differ by~64
from the codes of |?|, |@|, |A|; this convention applies only to
characters with ASCII codes less than~128. There are 256 possible characters,
hence 256 entries in each of the |\catcode|, |\mathcode|,
|\lccode|, |\uccode|, |\sfcode|, and |\delcode| tables. All
|\lccode|, |\uccode|, and |\char| values
must be less than~256. A font has at most 256 characters. There are
256~|\box| registers, 256~|\count| registers, 256~|\dimen| registers,
256~|\skip| registers, 256~|\muskip| registers, 256~|\toks| registers,
256~hyphenation tables.
The ``at size'' of a font must be less than~$2048\pt$, i.e.,~$2^{11}\pt$.
Math delimiters are encoded by multiplying the math~code of the ``small
character'' by~$2^{12}$. The magnitude of
a~\<dimen> value must be less than~$16384\pt$, i.e.,~$2^{14}\pt$;
similarly, the \<factor> in a~\<fil dimen> must be less than~$2^{14}$.
A~|\mathchar| or |\spacefactor| or |\sfcode| value must be less than~$2^{15}$;
a~|\mathcode| or |\mag| value must be less than or equal to~$2^{15}$,
and $2^{15}$ denotes an ``active'' math character. There
are $2^{16}\rm\,sp$ per~pt. A~|\delcode| value
must be less than~$2^{24}$; a~|\delimiter|, less than $2^{27}$.
The |\end| command sometimes contributes
a penalty of $-2^{30}$ to the current page. A~\<dimen> must be less than
$2^{30}\rm\,sp$ in absolute value; a~\<number> must be
less than $2^{31}$ in absolute value.

\endchapter

Mathematics is known in the trade as {\rm difficult,} or {\rm penalty, copy}
because it is slower, more difficult, and more expensive to set in type
than any other kind of copy normally occurring in books and journals.
\author UNIVERSITY OF ^{CHICAGO} PRESS, {\sl A Manual of Style\/} %
  (1969) % 12th edition, page 295

\bigskip

The tale of Math is a complex one,
and it resists both a simple plot summary
and a concise statement of its meaning.
\author PATRICK K. ^{FORD}, {\sl The Mabinogi\/} (1977) % p89
   % from his introduction to "Math Son of Mathonwy"

\eject
